\documentclass[11pt]{article}
\usepackage{amsmath, amssymb}
\usepackage{amsthm}
\usepackage{graphicx}
\usepackage{booktabs}
\usepackage{hyperref}
\usepackage[margin=1in]{geometry}
\newtheorem{conjecture}{Conjecture}
\title{On the Distribution of Smaller Factors in Semiprimes}
\author{Jonathan Kendall}
\date{April 24, 2026}

\begin{document}

\maketitle

\begin{abstract}
We study the distribution of semiprimes \(n = pq\) (with \(p \leq q\)) through the
normalized logarithmic coordinate \(\theta = \frac{\log p}{\log n}\).  Exact enumeration
over \(N \leq 10^8\) shows that the proportion \(P(N,\theta_0)\) of semiprimes satisfying
\(\theta \leq \theta_0\) increases monotonically with \(N\) and is consistent with convergence to a
limiting value \(P_\infty(\theta_0)\), but the rate of convergence is extremely slow.
The deficit \(\delta(N,\theta_0) = P_\infty(\theta_0) - P(N,\theta_0)\) remains
substantial even at \(N = 10^8\), consistent with \(O(1/\!\log\log N)\) decay arising
from the \(\log\log N\) accumulation of small prime contributions
(Montgomery--Vaughan \cite{montgomery2007}, \S7.4).
A second-order heuristic applied to the counting sum
$\sum_{p \leq N^{\theta_0}} \pi(N/p)$ (Section~\ref{sec:corrected}) suggests the
refined prefactor
\[
  f(\theta_0) = \log\!\left(\tfrac{1}{\theta_0}\right) - \theta_0 - \log 2 + \tfrac{1}{2},
\]
which is consistent with reducing the relative prediction error for
$b_{\mathrm{eff}} = \delta\cdot\log\log N$ from up to \(4.3\times\) (naive Mertens)
to under \(1.21\times\) across eight thresholds \(\theta_0\in\{0.10,\ldots,0.45\}\);
whether $C = \tfrac{1}{2}$ is the exact asymptotic constant remains open
(Conjecture~\ref{conj:prefactor}).
Over the observed range \(N \in [10^4, 10^8]\), \(\log\log N\) varies only from
\(2.22\) to \(2.91\), a range too narrow to distinguish \(1/\!\log\log N\) decay from
a local linear fit in \(\log N\).  Accordingly, the deficit is well approximated
\emph{locally} by
\[
  \delta(N,\theta_0) \;\approx\; a(\theta_0) - b(\theta_0)\,\log N
  \qquad (R^2 > 0.97 \text{ within } [10^4,10^8]),
\]
where \(b(\theta_0)\) is an empirical rate summarizing the behavior in this window.
The empirical regression \(b(\theta_0) \approx k(1-\theta_0)^{\alpha}\),
\(k \approx 0.015\), \(\alpha \approx 3.1\), is a descriptive fit to 4 data points.

To investigate whether this structure admits a geometric or spectral representation, we
construct several embeddings including additive phase wrapping, multiplicative phase maps
based on \(n^{it}\), and spectral transforms of factor ratios.  Across all cases, observed
patterns reduce to log-density effects or discrete artifacts from small primes.

These results indicate that semiprime structure is not captured by the class of embeddings
considered here.  The \(\theta\) coordinate provides the most stable description of factor
asymmetry, and the observed slow convergence is consistent with ratio-constraint phenomena
in multiplicative number theory.
\end{abstract}

\section{Introduction}

We study semiprimes \(n = pq\) with \(p \leq q\) and focus on the normalized coordinate
\[
  \theta = \frac{\log p}{\log n}.
\]
This parameter captures the relative size of the smaller factor and provides a natural
measure of factor asymmetry.

Classical results in analytic number theory describe the distribution of smooth numbers
and the frequency of integers with a smallest prime factor in a given range
(Dickman \cite{dickman1930}, Hildebrand--Tenenbaum \cite{hildebrand1993},
Tenenbaum \cite{tenenbaum1995}).  These results characterize the \emph{absolute} size
of the smallest prime factor relative to the integer, typically through the Dickman
function \(\rho(u)\) and its analogues.  They do not, however, explicitly study the
\emph{normalized} coordinate \(\theta = \log p / \log n\) for semiprimes, nor the
convergence of its distribution as \(N \to \infty\).  The present work provides
empirical evidence for this distribution and its convergence behavior, complementing
the theoretical framework with quantitative data over \(N \leq 10^8\).

The main contributions of this note are: (i) to document empirically the extremely slow
convergence of \(P(N,\theta_0)\) to its limit, characterizing the deficit locally as a
linear function of \(\log N\) over the range \(N \in [10^4, 10^8]\); and (ii) to show
that several natural geometric and spectral embeddings fail to recover this structure.

\section{Definitions}

Let \(n = pq\) be a semiprime with \(p \leq q\).  We define the normalized logarithmic
coordinate
\[
  \theta = \frac{\log p}{\log n}.
\]
For a fixed threshold \(\theta_0 < \frac{1}{2}\), define
\[
  P(N,\theta_0) = \Pr(\theta \leq \theta_0)
\]
over semiprimes \(n \leq N\).

The condition \(\theta \leq \theta_0\) is equivalent to
\[
  p \leq n^{\theta_0} = (pq)^{\theta_0},
\]
which implies
\[
  p^{1-\theta_0} \leq q^{\theta_0} \quad \Longrightarrow \quad
  p \leq q^{\,\theta_0/(1-\theta_0)}.
\]
Thus, \(\theta\) imposes a relative size constraint between the two prime factors rather
than a classical smoothness condition.

The algebraic identity
\begin{equation}\label{eq:identity}
  \log\!\left(\frac{p}{q}\right) = (2\theta - 1)\,\log n
\end{equation}
follows immediately: since \(\log p = \theta\log n\) and \(\log q = (1-\theta)\log n\),
one has \(\log(p/q) = \theta\log n - (1-\theta)\log n = (2\theta-1)\log n\).  This
identity holds exactly for every semiprime (Figure~\ref{fig:scatter}).

\section{Empirical Observations}

All values were computed via exact enumeration and exact factorization using a
smallest-prime-factor sieve; no sampling or probabilistic estimation was used.

\subsection{Convergence table}

Table~\ref{tab:convergence} records \(P(N,\theta_0)\) for
\(\theta_0 \in \{0.20,\,0.25,\,0.30,\,0.35\}\ \) and
\(N \in \{10^4,\,10^5,\,10^6,\,10^7,\,10^8\}\).

\begin{table}[h]
\centering
\caption{Proportion \(P(N,\theta_0)\) of semiprimes \(n \leq N\) satisfying
  \(\theta \leq \theta_0\).  All values in percent.}
\label{tab:convergence}
\begin{tabular}{rcccc}
\toprule
\(N\) & \(\theta_0 = 0.20\) & \(\theta_0 = 0.25\) & \(\theta_0 = 0.30\) & \(\theta_0 = 0.35\) \\
\midrule
\(10^4\) & 49.52 & 59.28 & 67.73 & 76.46 \\
\(10^5\) & 51.93 & 60.77 & 69.40 & 77.27 \\
\(10^6\) & 53.55 & 62.47 & 70.51 & 78.39 \\
\(10^7\) & 55.53 & 63.74 & 71.83 & 79.19 \\
\(10^8\) & 56.29 & 64.95 & 72.74 & 79.95 \\
\bottomrule
\end{tabular}
\end{table}

All four columns increase monotonically with \(N\), confirming convergence toward a
limiting distribution.  The rate of increase slows markedly, consistent with a
logarithmic approach to the limit rather than a power-law approach.

\subsection{Theta and factor ratio}

Figure~\ref{fig:scatter} shows the joint distribution of \(\theta\) and
\(\log(p/q)\) for all semiprimes \(n \leq 10^6\).  The scatter lies exactly on the
one-parameter family of lines \(\log(p/q) = (2\theta-1)\log n\) (identity
\eqref{eq:identity}), confirming that \(\theta\) fully determines the factor ratio once
\(n\) is fixed.  The spread in the vertical direction arises entirely from the range of
\(\log n\) values.

\begin{figure}[h]
\centering
\includegraphics[width=0.82\textwidth]{theta_vs_logratio.png}
\caption{Factor-ratio $\log(p/q)$ versus $\theta$ for all semiprimes $n \leq 10^6$
  (30\,000 sampled points shown; colour encodes $\log n$).
  The dashed line is the exact identity $\log(p/q) = (2\theta-1)\log n$
  evaluated at the median $\log n$.}
\label{fig:scatter}
\end{figure}

\subsection{Local linear approximation to the deficit}

Define the deficit \(\delta(N,\theta_0) = 1 - P(N,\theta_0)\).  Table~\ref{tab:convergence}
shows that \(\delta\) decreases with \(N\), but very slowly: at \(\theta_0 = 0.25\) the
deficit falls only from \(40.72\%\) to \(35.05\%\) as \(N\) grows by four orders of
magnitude.

\paragraph{Theoretical rate.}
The asymptotic rate of convergence is \(O(1/\!\log\log N)\), arising from the
\(\log\log N\) accumulation of small prime contributions in semiprime counting
(Montgomery--Vaughan \cite{montgomery2007}, \S7.4).  Over our observed range,
\(\log\log N\) increases from \(2.22\) (at \(N = 10^4\)) to only \(2.91\) (at
\(N = 10^8\)) — a total variation of \(0.69\) across four decades.  This near-constancy
means that \(1/\!\log\log N\) behaves almost linearly in \(\log N\) within this window,
making the two functional forms indistinguishable from the data alone.

\paragraph{Local linear fit.}
Within the observed range, a linear regression in \(\log N\) therefore provides an
accurate \emph{local} description:
\begin{equation}\label{eq:model}
  \delta(N,\theta_0) \;\approx\; a(\theta_0) - b(\theta_0)\,\log N
  \qquad \text{(empirical fit, } N \in [10^4, 10^8]\text{)}.
\end{equation}
This is not an asymptotic law: extrapolating \eqref{eq:model} to large \(N\) would
predict \(\delta < 0\), which is impossible.  It is a finite-range proxy for the
slower \(O(1/\!\log\log N)\) decay.  Numerical comparison confirms that the linear fit
(mean \(R^2 = 0.989\) across four thresholds) outperforms a direct \(c/\!\log\log N\)
fit (mean \(R^2 = 0.41\)) precisely because the variation of \(\log\log N\) is too
small in this window to distinguish the two forms; the linear model has more
flexibility with its two free parameters.

Table~\ref{tab:fit} records the regression parameters for each \(\theta_0\).

\begin{table}[h]
\centering
\caption{Parameters of the local linear fit \(\delta \approx a(\theta_0) - b(\theta_0)\log N\)
  over \(N \in [10^4, 10^8]\) (\(R^2 > 0.97\) in all cases).  This fit is a
  finite-range proxy for the theoretically predicted \(O(1/\!\log\log N)\) decay;
  it is not valid outside the observed range.}
\label{tab:fit}
\begin{tabular}{cccr}
\toprule
\(\theta_0\) & \(a(\theta_0)\) & \(b(\theta_0)\) & \(R^2\) \\
\midrule
0.20 & 0.5692 & 0.007444 & 0.9754 \\
0.25 & 0.4635 & 0.006216 & 0.9955 \\
0.30 & 0.3702 & 0.005402 & 0.9904 \\
0.35 & 0.2709 & 0.003866 & 0.9954 \\
\bottomrule
\end{tabular}
\end{table}

The rate \(b(\theta_0)\) varies with \(\theta_0\) and is analyzed further in
Section~\ref{sec:bscaling}.

\subsection{Scaling of the convergence rate \texorpdfstring{$b(\theta_0)$}{b(theta)}}
\label{sec:bscaling}

The empirical rate \(b(\theta_0)\) from Table~\ref{tab:fit} decreases as \(\theta_0\)
increases, meaning that a larger imbalance threshold corresponds to a faster-decreasing
deficit within the observed range.  We fit two parametric curves to the four data points
as \emph{descriptive regressions}, not theoretical predictions:

\paragraph{Model~1 (empirical).}
\[
  b(\theta_0) \approx k\,(1-\theta_0)^{\alpha},
  \qquad k \approx 0.01504,\quad \alpha \approx 3.056,\quad R^2 = 0.974.
\]

\paragraph{Model~2 (empirical).}
\[
  b(\theta_0) \approx \frac{c}{\theta_0^{\,\gamma}},
  \qquad c \approx 1.300\times 10^{-3},\quad \gamma \approx 1.109,\quad R^2 = 0.938.
\]

Model~1 fits the four data points more closely (\(R^2 = 0.974\) vs.\ \(0.938\)).
Both forms predict \(b \to 0\) as \(\theta_0 \to \tfrac{1}{2}\), qualitatively
consistent with the known scarcity of nearly balanced semiprimes.
\emph{Caution}: these regressions are fit to only four values of \(\theta_0\) and
should be understood as summaries of the observed behavior, not as laws with theoretical
backing.  Extending the table to additional \(\theta_0\) values and larger \(N\) would
be needed to discriminate between functional forms.

\subsection{Refined heuristic for the prefactor}
\label{sec:corrected}

The naive Mertens argument gives a leading prefactor of $\log(1/\theta_0)$ for the
$O(1/\!\log\log N)$ deficit.  We develop a \emph{second-order heuristic} — not a
proven asymptotic — that accounts for the next-order term in the prime sum and is
consistent with the data substantially better across the range $\theta_0\in[0.10,0.45]$.

\paragraph{Expansion of $1/\log(N/p)$.}
Applying $\pi(x)\approx x/\log x$ to $L(N,\theta_0)$ yields terms $N/(p\log(N/p))$.
Writing $\log(N/p)=\log N - \log p$ and expanding:
\[
  \frac{1}{\log(N/p)}
  = \frac{1}{\log N}\cdot\frac{1}{1-\tfrac{\log p}{\log N}}
  = \frac{1}{\log N}\!\left(1 + \frac{\log p}{\log N}
    + O\!\left(\frac{\log^2\!p}{\log^2\!N}\right)\right).
\]
Summing the first correction term over primes via Chebyshev's theorem,
$\sum_{p\leq x}(\log p)/p = \log x + O(1)$, produces an additive shift of
$(\log x)/\log N$ in each counting sum, evaluated at its respective cutoff.

\paragraph{Heuristic sums.}
Define
\[
  T(N) = \!\sum_{p\leq\sqrt{N}}\!\bigl(\pi(N/p) - \pi(p-1)\bigr), \qquad
  L(N,\theta_0) = \!\sum_{p\leq N^{\theta_0}}\!\bigl(\pi(N/p) - \pi(p-1)\bigr).
\]
Retaining the first correction gives the heuristic approximations
\begin{align*}
  T(N) &\;\approx\; \frac{N}{\log N}\!\left[\log\log N - \log 2 + M + \tfrac{1}{2}\right], \\[4pt]
  L(N,\theta_0) &\;\approx\; \frac{N}{\log N}\!\left[\log\log N + \log\theta_0 + M + \theta_0\right],
\end{align*}
where $M\approx 0.2615$ is the Meissel--Mertens constant.  The terms $\tfrac{1}{2}$
and $\theta_0$ arise from Chebyshev's theorem at the cutoffs $\sqrt{N}$ and $N^{\theta_0}$
respectively.  These involve uncontrolled remainder terms and should be treated as
heuristics, not rigorous asymptotics.

\paragraph{Refined prefactor.}
Forming $\delta = 1 - L/T$ from these heuristic sums gives the deficit as
\emph{well approximated by}
\begin{equation}\label{eq:prefactor}
  \delta(N,\theta_0) \;\approx\;
  \frac{f(\theta_0)}{\log\log N + c},
\end{equation}
where $c = -\log 2 + M + \tfrac{1}{2} \approx 0.07$ and
\begin{equation}\label{eq:f}
  f(\theta_0) = \log\!\left(\tfrac{1}{\theta_0}\right) - \theta_0 - \log 2 + \tfrac{1}{2}.
\end{equation}
We do not claim \eqref{eq:f} is the proven asymptotic prefactor; it is the output of a
second-order heuristic that is consistent with the data to within 20\% across the
range tested.  The residual 20\% gap motivates the following formal statement.

\begin{conjecture}\label{conj:prefactor}
The limit $\lim_{N\to\infty}\delta(N,\theta_0)\cdot\log\log N$ exists and equals a
function $f(\theta_0)$ satisfying
\[
  f(\theta_0) = \log\!\left(\tfrac{1}{\theta_0}\right) - \theta_0 - \log 2 + C
\]
for a universal constant $C$ (independent of $\theta_0$).  Empirically, $C \approx
\tfrac{1}{2}$ is consistent with the data across $\theta_0 \in [0.10,\,0.45]$.
\end{conjecture}

The value $C = \tfrac{1}{2}$ arises in the heuristic from Chebyshev's sum at the
cutoff $\sqrt{N}$ of $T(N)$; whether it is exact in the asymptotic sense is not
established by the present data.

\paragraph{Numerical comparison.}
Table~\ref{tab:prefactor} compares the naive prefactor $\log(1/\theta_0)$, the
heuristic $f(\theta_0)$, and the empirical $b_{\mathrm{eff}} = \delta\cdot\log\log N$
averaged over $N\in[10^4,10^8]$.

\begin{table}[h]
\centering
\caption{Empirical $b_{\mathrm{eff}}$ (mean of $\delta\cdot\log\log N$ over
  $N\in[10^4,10^8]$) vs.\ the naive prefactor $\log(1/\theta_0)$ and the heuristic
  $f(\theta_0)=\log(1/\theta_0)-\theta_0-\log 2+\tfrac{1}{2}$ (Conjecture~\ref{conj:prefactor}).}
\label{tab:prefactor}
\begin{tabular}{cccccc}
\toprule
$\theta_0$ & $\log(1/\theta_0)$ & $f(\theta_0)$ & $b_{\mathrm{eff}}$
           & $b_{\mathrm{eff}}/\log(1/\theta_0)$ & $b_{\mathrm{eff}}/f(\theta_0)$ \\
\midrule
0.10 & 2.303 & 2.009 & 1.819 & 0.790 & 0.905 \\
0.15 & 1.897 & 1.554 & 1.465 & 0.772 & 0.943 \\
0.20 & 1.609 & 1.216 & 1.205 & 0.749 & 0.991 \\
0.25 & 1.386 & 0.943 & 0.976 & 0.704 & 1.034 \\
0.30 & 1.204 & 0.711 & 0.763 & 0.634 & 1.074 \\
0.35 & 1.050 & 0.507 & 0.562 & 0.535 & 1.108 \\
0.40 & 0.916 & 0.323 & 0.373 & 0.407 & 1.154 \\
0.45 & 0.799 & 0.155 & 0.187 & 0.234 & 1.202 \\
\bottomrule
\end{tabular}
\end{table}

The heuristic is consistent with reducing the maximum relative error from $4.3\times$
(naive) to under $1.21\times$ across all eight thresholds.  The remaining monotone
trend in $b_{\mathrm{eff}}/f(\theta_0)$ (from $0.90$ to $1.20$) suggests further
subleading structure, partly addressed in the next paragraph.

\paragraph{Discrete activation thresholds (empirical observation).}
A separately observable feature of the convergence is the presence of discrete steps
in $P(N,\theta_0)$.  A prime $p$ begins contributing to $L(N,\theta_0)$ only at
$N = p^{1/\theta_0}$, at which point the number of valid lopsided pairs $(p,q)$ with
$n = pq \leq N$ increases by a finite amount.  These \emph{activation thresholds}
produce step-like upward jumps in $L$ (and hence in $P$) at each $N = p^{1/\theta_0}$.

For small $\theta_0$ the thresholds are widely spaced and large.  At $\theta_0=0.10$,
the primes $2,3,5,7$ activate at $N = 2^{10} = 1\,024$,\; $3^{10}\approx 59\,000$,\;
$5^{10}\approx 9.8\times 10^6$,\; $7^{10}\approx 2.8\times 10^8$.  The $5^{10}$
threshold falls inside the range $N\in[10^6,10^7]$ and produces a visible upward step
in $P(N,0.10)$ between those values, creating the non-monotone behaviour in
$b_{\mathrm{eff}}$ at $\theta_0\in\{0.10,0.15\}$ visible in Table~\ref{tab:prefactor}.

This discrete-threshold mechanism is a genuinely novel empirical observation: the
convergence of $P(N,\theta_0)$ is not smooth in $N$ but proceeds in steps indexed
by the prime activation sequence $\{p^{1/\theta_0}\}$.  It also explains why
smooth models for the deficit perform less well at small $\theta_0$ over this range.

\section{Geometric and Spectral Tests}

We tested several candidate embeddings:
\begin{itemize}
  \item Additive phase wrapping: \(\phi_t(n) = t\log n \bmod 2\pi\)
  \item Multiplicative phase maps using \(n^{it}\)
  \item Spectral fingerprint \(C(t) = \bigl|\sum e^{it\log n}\bigr|\,/\,N\)
  \item Factor-ratio spectrum \(C_{\mathrm{ratio}}(t) = \bigl|\sum e^{it\log(p/q)}\bigr|\,/\,N\)
\end{itemize}

Localization was assessed by examining stability under variation of \(N\) and clustering
by \(\theta\).  No stable localization was observed.

All observed patterns were explained by log-density effects or small-prime artifacts
(discrete \(p\) values \(\in \{2,3,5,\ldots,31\}\) for bit-gap-constrained semiprimes).
None of the embeddings captured factor asymmetry.

\section{Results}

\subsection{Summary of empirical findings}

For fixed \(\theta_0 \in (0, \tfrac{1}{2})\), the following is established empirically:
\begin{enumerate}
  \item \(P(N,\theta_0)\) increases monotonically with \(N\) for all four tested thresholds.
  \item The increase is extremely slow: over four decades of \(N\), the gain is at most
    \(6.8\%\) (at \(\theta_0 = 0.20\)), compared with the \(\sim 50\%\) gap that remains
    at \(N = 10^8\).
  \item Within the observed range, the deficit \(\delta(N,\theta_0) = 1 - P(N,\theta_0)\)
    decreases at a rate well approximated by a constant per unit \(\log N\), with rate
    \(b(\theta_0)\) as tabulated in Table~\ref{tab:fit}.
\end{enumerate}

\subsection{Interpretation and connection to known theory}

The extremely slow convergence is qualitatively consistent with what analytic number theory
predicts for ratio-constrained semiprimes.  The density of semiprimes \(n = pq \leq N\)
with \(p \leq N^{u}\) for fixed \(u < \tfrac{1}{2}\) is governed by sums of the form
\(\sum_{p \leq N^u} \pi(N/p)\), where \(\pi\) is the prime counting function.  Standard
sieve estimates (see Hildebrand--Tenenbaum \cite{hildebrand1993} and
Tenenbaum \cite{tenenbaum1995}) show that such sums accumulate logarithmic corrections
involving factors of \(\log\log N / \log N\), producing convergence that is invisible
over a few decades but genuine asymptotically.

The local linear fit \eqref{eq:model} captures the dominant behavior in the observed
range \([10^4, 10^8]\), but it is a finite-range approximation: the true deficit decays
to zero more slowly than any fixed linear function of \(\log N\) would predict, and
the functional form of \(P_\infty(\theta_0)\) is not determined by our data.  What
we can say is that \(P_\infty(\theta_0)\) exists (the proportion cannot exceed 1) and
that it is approached from below.

\subsection{The \(\theta\) coordinate}

The exact identity \(\log(p/q) = (2\theta-1)\log n\) (equation~\eqref{eq:identity})
shows that \(\theta\) and \(\log n\) together determine the factor ratio with no residual
freedom.  This makes \(\theta\) a natural coordinate: it separates the role of the
factor imbalance (encoded in \(\theta\)) from the role of the magnitude of \(n\)
(encoded in \(\log n\)).  The observed convergence is therefore a statement about how
the distribution of \(\theta\) across semiprimes \(n \leq N\) evolves as \(N\) grows,
which is a well-posed and nontrivial question even if its full resolution requires
deeper sieve-theoretic tools.

\section{Conclusion}

Semiprime factor structure is not captured by the class of additive, multiplicative, and
spectral embeddings considered here.  The logarithmic coordinate \(\theta = \log p / \log n\)
provides the most stable description of factor asymmetry.

The rate of convergence is consistent with \(O(1/\!\log\log N)\), as predicted by the
\(\log\log N\) accumulation of small prime contributions in semiprime counting
(Montgomery--Vaughan \cite{montgomery2007}, \S7.4).  The local linear-in-\(\log N\)
approximation of Section~3.3 is a valid finite-range description of this slower decay:
over \(N \in [10^4, 10^8]\), \(\log\log N\) varies by only \(0.69\), so the two
functional forms are numerically indistinguishable in this window.

The novel contributions are: (1)~empirical quantification of the convergence rate
$b(\theta_0)$ across four imbalance thresholds, suggesting it decreases smoothly with
$\theta_0$; (2)~the geometric and spectral null results, consistent with the
$\theta$ coordinate being unrecoverable from the embeddings tested; (3)~a second-order
heuristic for the prefactor $f(\theta_0) = \log(1/\theta_0) - \theta_0 - \log 2 +
\tfrac{1}{2}$ (Conjecture~\ref{conj:prefactor}, Section~\ref{sec:corrected}),
consistent with reducing the naive-Mertens prediction error from $4.3\times$ to
under $1.21\times$ by accounting for the asymmetric Chebyshev sums at the cutoffs
$\sqrt{N}$ and $N^{\theta_0}$; and (4)~the empirical observation that $P(N,\theta_0)$
increases in discrete steps at $N = p^{1/\theta_0}$ for each prime $p$ (activation
thresholds), which is a genuinely novel structural feature of the convergence not
captured by smooth asymptotic models.
Proving Conjecture~\ref{conj:prefactor} and determining $P_\infty(\theta_0)$
rigorously remain open problems accessible to sieve-theoretic methods.

The negative results on geometric and spectral embeddings suggest that the logarithmic
coordinate is not merely convenient but is in some sense the correct one for this problem.

\begin{thebibliography}{9}

\bibitem{dickman1930}
K.~Dickman,
``On the frequency of numbers containing prime factors of a certain relative
  magnitude,''
\textit{Arkiv f\"{o}r Matematik, Astronomi och Fysik},
vol.~22A, no.~10, 1930.

\bibitem{hildebrand1993}
A.~Hildebrand and G.~Tenenbaum,
``Integers without large prime factors,''
\textit{Journal de Th\'{e}orie des Nombres de Bordeaux},
vol.~5, pp.~411--484, 1993.

\bibitem{tenenbaum1995}
G.~Tenenbaum,
\textit{Introduction to Analytic and Probabilistic Number Theory},
Cambridge University Press, 1995.

\bibitem{erdos1979}
P.~Erd\H{o}s and C.~Pomerance,
``On the distribution of smooth numbers,''
in \textit{Number Theory, Carbondale 1979}, Lecture Notes in Math.\ 751,
Springer, Berlin, 1979, pp.~173--183.

\bibitem{halberstam1974}
H.~Halberstam and H.-E.~Richert,
\textit{Sieve Methods},
Academic Press, London, 1974.

\bibitem{montgomery2007}
H.~L.~Montgomery and R.~C.~Vaughan,
\textit{Multiplicative Number Theory~I: Classical Theory},
Cambridge Studies in Advanced Mathematics~97,
Cambridge University Press, 2007.

\end{thebibliography}

\end{document}
